\part{XÂY DỰNG ĐỒ ÁN TỔNG HỢP}

\chapter{Phân tích và Thiết kế hệ thống}

\section{Mô tả bài toán}
Giải quyết vấn đề quản lý công việc và theo dõi tiến độ trong các nhóm làm việc (Todo $\rightarrow$ Doing $\rightarrow$ Done). Hệ thống cho phép người dùng trực quan hóa quy trình làm việc thông qua bảng Kanban.

\section{Chức năng chính}
\begin{itemize}
    \item \textbf{Authentication:} Đăng nhập/Đăng xuất đơn giản với Protected Routes (Chỉ những người đã đăng nhập mới xem được Board).
    \item \textbf{Board Dashboard:} Trang danh sách các dự án (Board).
    \item \textbf{Kanban Board:} Gồm nhiều danh sách (List).
    \item \textbf{Card Management:} Thêm, sửa, xóa các thẻ công việc (Card) trong List.
    \item \textbf{Drag and Drop:} Kéo thả card giữa các list để thay đổi trạng thái công việc.
\end{itemize}

\section{Kiến trúc linh kiện (Component Tree)}
\begin{verbatim}
App
 +- RouterProvider
    +- AuthProvider
       +- BoardProvider
          +- MainLayout
             +- Pages (Login, Dashboard, BoardDetail)
                +- Features/Boards/Components
                   +- BoardCard
                   +- List
                      +- TaskCard
                      +- AddCardForm
                   +- AddListForm
                   +- CardModal
\end{verbatim}

\chapter{Triển khai và Demo sản phẩm}

\section{Tích hợp hệ thống và Luồng dữ liệu (State Flow)}
Khi người dùng thêm một Card mới:
\begin{enumerate}
    \item Nhập tiêu đề vào \texttt{AddCardForm}.
    \item Dispatch hành động \texttt{ADD\_CARD} thông qua \texttt{useCards}.
    \item \texttt{boardReducer} cập nhật trạng thái trong Context.
    \item React tự động re-render các \texttt{List} bị thay đổi dữ liệu.
\end{enumerate}

\section{Tính năng nâng cao}
\begin{itemize}
    \item \textbf{Local Storage Persistence:} Lưu trữ dữ liệu trực tiếp trong trình duyệt, thông tin không bị mất khi làm mới trang.
    \item \textbf{Card Detail Modal:} Popup chi tiết công việc cho phép quản lý nhãn (Labels), ngày hết hạn (Due Date) và mô tả dài.
    \item \textbf{Dynamic Form Inputs:} Sử dụng các form inline (\texttt{AddCardForm}, \texttt{AddListForm}) để tăng tính tương tác thay vì chuyển trang.
\end{itemize}

\section{Hướng dẫn cài đặt (Deployment)}
\begin{enumerate}
    \item Giải nén mã nguồn.
    \item Chạy \texttt{npm install} để cài đặt các thư viện cần thiết (\texttt{react-router-dom}, \texttt{dnd-kit}).
    \item Chạy \texttt{npm run dev} để khởi động môi trường phát triển.
\end{enumerate}
